% JAC1964-20261121. /root/programme_audit; GPT-6 Astra Pro.
% Accepted input: 032 Corollary 8.3/Milne's full Tate consequence.
% Additional established input: Kahn 2003 Theorem 1.10.

\subsection*{JAC-P15. The numerical--tau assertion over finite fields}

\begin{theorem}[JAC-P15]\label{jac:P15}
Accept the Tate consequence of 032, Corollary~8.3, printed
pp.~49--50, as also formulated in the supplied Milne report,
Theorem~3.1, printed p.~6. Let $k$ be a finite field
or an algebraic closure of a finite field. Let $X/k$ be smooth and
projective, and suppose that its geometric rational Chow motive is
of abelian type: it belongs to the thick rigid tensor category
generated by abelian varieties and Artin motives.
Then, for every $r$, the natural map
\[
 \operatorname{CH}^r(X)_{\mathbf Q}
        \longrightarrow N^r(X)_{\mathbf Q}                 \tag{JAC.P15.1}
\]
is an isomorphism, where $N^r$ denotes cycles modulo numerical
equivalence. In particular, for an integral cycle $z$ on $X$,
\[
 z\equiv_{\rm num}0
 \quad\Longleftrightarrow\quad
 [z]\text{ is torsion in }\operatorname{CH}^r(X)
 \quad\Longleftrightarrow\quad
 z\equiv_\tau0.                                           \tag{JAC.P15.2}
\]
Here $\tau$ has the customary meaning that a nonzero multiple is
algebraically equivalent to zero.

Consequently cote~44's Th.~7, with that meaning of $\tau$, holds
for every abelian variety over $\overline{\mathbf F}_p$, in every
dimension and codimension. The conclusion also holds for smooth
projective varieties whose geometric Chow motive is of abelian
type, including products of curves.
\end{theorem}

\begin{proof}
We separate the accepted arithmetic input from the additional
established motivic theorem.

\emph{Tate passes to abelian-type summands.}
The accepted 032, Corollary~8.3, printed pp.~49--50, says that for every abelian variety
$A/\mathbf F_q$, every $r$ and every $\ell\ne p$, the cycle-class
map
\[
 \operatorname{CH}^r(A)\otimes\mathbf Q_\ell
   \longrightarrow H^{2r}(A_{\overline{\mathbf F}_q},
                         \mathbf Q_\ell(r))^{G_q}
\]
is surjective. This applies to every product of abelian varieties
and after every finite extension of the finite field.

Every construction of an object of the thick rigid category uses
finitely many sums, tensor products, duals, projectors, and splitting
correspondences. After a finite extension $k'/k$, its Artin factors
split. Tensor products of abelian motives are motives of abelian
products; their duals are Tate twists of such motives, by Poincare
duality. Thus, after enlarging $k'$ if necessary, $h(X_{k'})$ is
a direct summand of a finite direct sum
\[
              \bigoplus_j h(A_j)(t_j)                    \tag{JAC.P15.3}
\]
with each $A_j/k'$ abelian. If the initial abelian-type presentation
was given only over the algebraic closure, the same conclusion
holds: the varieties, projectors, correspondences and the finitely
many rational equivalences giving their splitting identities all
descend to a finite field extension.

A Galois-invariant cohomology class on $X_{k'}$ maps, by the
inclusion in (JAC.P15.3), to invariant classes of the appropriate
degrees and Tate twists on the $A_j$. The accepted Tate statement
represents them by $\mathbf Q_\ell$-linear combinations of cycles.
Applying the algebraic retraction in (JAC.P15.3) gives a cycle
representing the original class on $X_{k'}$. Thus $X_{k'}$ satisfies
Tate in every codimension. If the class originally came from a
$G_k$-invariant class on $X$, its cycle representative descends
rationally by the normalized norm
\[
           [k':k]^{-1}\pi_*,\qquad \pi:X_{k'}\to X.
\]
Indeed pullback of this norm is the Galois average and fixes that
class. No assertion that every tensor product of previously known
Tate classes spans all tensor invariants is used; Tate is invoked
on the abelian product itself.

\emph{Kahn's theorem upgrades the equivalence relation.}
Kahn, \emph{Equivalences rationnelle et numerique sur certaines
varietes de type abelien sur un corps fini}, Ann.\ Sci.\ Ecole
Norm.\ Sup.\ (4) 36 (2003), 977--1002, Definition~1 and Theorem~1,
printed p.~978, and Theorem~1.10, printed p.~982, proves that
rational and numerical equivalence agree with rational
coefficients for an abelian-type variety satisfying Tate.
All of these hypotheses have just been verified for $X_{k'}$.
We therefore obtain (JAC.P15.1) over $k'$.

For precision, the mechanism in that source is stronger than
mere equality of homological and numerical equivalence. The
nilpotence theorem for Kimura motives lifts the numerical simple
decomposition. A non-unit simple factor on which Frobenius is
not $1$ has no morphism from the unit: if its minimal Frobenius
polynomial is $P$ with $P(1)\ne0$, nilpotence gives
$P(F)^a=0$, whereas $Ff=f$ for a morphism $f$ from the unit.
The unit factors lift to actual unit motives. Tate and Frobenius
semisimplicity identify exactly these factors. The precise source
loci are Theorem~1.4, Proposition~1.5, Lemma~1.6 and
Proposition~1.7, printed p.~981, and Lemma~1.9 and
Theorem~1.10, printed p.~982.

Suppose now that $k$ is finite and $z$ is numerically trivial on
$X$. Its pullback to $k'$ is numerically trivial: for each
complementary cycle $w$ on $X_{k'}$, the projection formula gives
\[
 [k':k]\,\deg_{k'}(\pi^*z\cdot w)
       =\deg_k(z\cdot\pi_*w)=0.
\]
Hence $\pi^*z=0$ in rational Chow, and
$\pi_*\pi^*z=[k':k]z$ gives $z=0$ there as well.
This proves (JAC.P15.1) even when the abelian-type presentation
only becomes available over $k'$.

If $k=\overline{\mathbf F}_p$, choose a finite field of definition
for $X$, $z$, and its abelian-type presentation. Numerical
triviality of $z$ over $k$ implies numerical triviality over that
finite field, since its test cycles form a subset of the geometric
test cycles. The finite-field conclusion makes $z$ zero in
rational Chow there, hence also over $k$. The reverse implication
in (JAC.P15.1) follows from rational equivalence implying numerical
equivalence. Surjectivity of (JAC.P15.1) is built into the
definition of the numerical quotient.

Finally, an element of an abelian group vanishes after tensoring
with $\mathbf Q$ exactly when it is torsion. Thus numerical
triviality implies Chow torsion. Chow torsion implies
$\tau$-triviality because a rationally trivial multiple is
algebraically trivial. Conversely a $\tau$-trivial cycle is
numerically trivial: all intersection numbers of its algebraically
trivial multiple vanish, and integer intersection numbers have
no nonzero torsion. This proves (JAC.P15.2).
\end{proof}

\paragraph{A finite-field consequence for higher $K$-theory.}
For $X$ as in JAC-P15 over a finite field,
\[
                    K_i(X)\otimes\mathbf Q=0
                    \qquad(i>0).                         \tag{JAC.P15.4}
\]
For $X$ already of abelian type over that field, this is Kahn's
Corollary~2.2, printed p.~982, applied to the Tate result above.
In the geometric abelian-type case, apply that corollary after
the finite extension giving (JAC.P15.3), then use the restriction
and transfer identity $\mathrm{cor}\,\mathrm{res}=[k':k]$ on
$K_i(X)$. This is a consequence through Kahn's cited theorem,
not a new proof of the general Parshin conjecture.

\paragraph{What changes in the programme.}
JAC-P06 correctly identifies the finite-field meaning of $\tau$,
but does not itself prove the numerical comparison. JAC-P15 now
supplies that comparison using the \emph{Tate} consequence of the
accepted input and Kahn's additional established theorem. It
does not follow from the Hodge standard statement of Corollary~8.4
alone. It does not apply to every smooth projective variety in
positive characteristic without the abelian-type hypothesis, to
arbitrary transcendental extensions of $\overline{\mathbf F}_p$,
or to the number-field Fermat Jacobian of JAC-P05. The latter
retains its non-torsion Ceresa obstruction.

\paragraph{Primary source verification.}
The published 26-page Kahn article was directly read from
\texttt{https://www.numdam.org/item/ASENS\_2003\_4\_36\_6\_977\_0.pdf}.
The author's 35-page \texttt{cell4.pdf} has the same theorem numbers
but different pagination: Definition~1 p.~2, Theorem~1 p.~3,
the nilpotence argument pp.~6--7, and Theorem~1.10 p.~8.
The older arXiv version uses different numbering and is not used
for these exact locators. The conclusion above is a deduction
in this run from accepted 032/Milne and established Kahn; no
historical originality or formal verification is asserted.
